% Abstracts replaced on 2026-11-06 (README (172)); kept verbatim, not compiled.

% ---- from main.tex ----
We study positive-data learning of multiple context-free languages under a
fixed finite-semigroup typing: a morphism $s$ from $\Sigma^+$ onto a finite
semigroup $\mathsf S$ is fixed in advance, and nonempty factors and tuple
components are compared only when their $s$-values agree.  This continues the
context-free case of a companion paper.  Identification in the limit follows
Yoshinaka's distributional framework; our subject is the total
length of characteristic samples, measured against the size $\|G\|$ and the
rule thickness $\bar\theta_G$ of a presentation $G$ of the target.  We show
that local groups give polynomial characteristic data at every fan-out and are
the exact pseudovariety boundary at fan-out one.  If
$\mathsf S$ is a local group, then for every fixed fan-out and rank bound, each
target has characteristic data of length $O(\|G\|^2\bar\theta_G)$ for each of
its reduced good presentations $G$ of bounded rank, so these classes are
identifiable in polynomial time and thick data.  The factor
$\|G\|\bar\theta_G$ is unavoidable: at fan-out $f$, over an alphabet of at
least $4f+1$ letters (two at fan-out one) and under every typing, some finite
targets with near-minimal presentations need data of length
$\Omega(\|G\|\bar\theta_G)$, which matches Yoshinaka's upper bound for the trivial
typing.  At fan-out one and rank at least two, local groups are exactly the
pseudovarieties of finite semigroups with uniformly polynomial characteristic
data, and hence exactly those whose classes are identifiable in polynomial time
and thick data; outside them, targets with presentations of size and rule
thickness $O(n)$ need data of length exponential in $n$, for every set-driven
reconstruction operator that has characteristic samples throughout the class.  This witness family does not carry over to fan-out two: under any
fixed typing, all but finitely many of its members fail to be
tuple-substitutable there, and no target there keeps their long word as the
only one of its type.  Whether local groups remain the boundary at fan-out two
is left open.  The upper bound rests on a simultaneous shortening of
interleaved tuple components that preserves ordered products in finite groups.

% ---- from japanese/main_JP.tex ----
本稿は，固定した有限半群型付けの下での多重文脈自由言語の正データ学習を研究する：
準同型 $s$ を，$\Sigma^+$ から有限半群 $\mathsf S$ の上への写像として前もって固定し，
空でない因子とタプル成分は，それらの $s$-値が一致するときにのみ比較する．
これは姉妹論文の文脈自由の場合の続きである．
極限同定は Yoshinaka の分布的枠組みに従う；本稿の主題は特徴標本の総長であり，
それを大きさ $\|G\|$ と規則の厚さ $\bar\theta_G$ に対して測る（これらは標的の表示 $G$ のものである）．
本稿は，局所群がすべてのファンアウトで多項式の特徴データを与え，
ファンアウト1ではちょうど擬多様体の境界であることを示す．
$\mathsf S$ が局所群ならば，固定したすべてのファンアウトとランクの上限について，
各標的は長さ $O(\|G\|^2\bar\theta_G)$ の特徴データを，ランク有界なその既約な良い表示 $G$ のそれぞれについてもち，
したがってこれらのクラスは多項式時間・厚データ同定可能である．
因子
$\|G\|\bar\theta_G$ は避けられない：ファンアウト $f$ において，
$4f+1$ 文字以上（ファンアウト1では2文字）のアルファベット上で，かつすべての型付けの下で，
最小に近い表示をもつある有限の標的は長さ
$\Omega(\|G\|\bar\theta_G)$ のデータを必要とし，これは自明な型付けに対する Yoshinaka の上界と一致する．
ファンアウト1かつランク2以上では，局所群は一様に多項式の特徴データをもつ有限半群の擬多様体にちょうど一致し，
したがってそのクラスが多項式時間・厚データ同定可能なものにちょうど一致する；
それらの外では，大きさと規則の厚さが $O(n)$ の表示をもつ標的が，$n$ について指数的な長さのデータを必要とする．
これは，クラス全体で特徴標本をもつすべての集合駆動型再構成作用素について成り立つ．
この証拠族はファンアウト2には持ち越されない：任意の固定した型付けの下で，
その成員のうち有限個を除くすべてはそこでタプル代入可能でなく，
そこでのどの標的も，それらの長い語をその型の唯一の語として保たない．
局所群がファンアウト2でも境界であり続けるかどうかは未解決のまま残す．
上界は，有限群における順序付き積を保つ，交互に並ぶタプル成分の同時短縮に基づく．
